\pagebreak
\section{Entity and Relationship Extraction Approach}
\label{app:prompts}

The following prompts, designed for GPT-4, are used in the default GraphRAG initialization pipeline:
\begin{itemize}
\item \href{https://github.com/microsoft/graphrag/blob/6d21ef268377e319a165ca2250bd6841737df1ad/graphrag/prompts/index/entity_extraction.py#L6}{Default Graph Extraction Prompt}
\item \href{https://github.com/microsoft/graphrag/blob/6d21ef268377e319a165ca2250bd6841737df1ad/graphrag/prompts/index/claim_extraction.py#L1}{Claim Extraction Prompt}
\end{itemize}

\subsection{Entity Extraction}
\label{app:entity}
We do this using a multipart LLM prompt that first identifies all \emph{entities} in the text, including their name, type, and description, before identifying all \emph{relationships} between clearly related entities, including the source and target entities and a description of their relationship. Both kinds of element instance are output in a single list of delimited tuples.

\subsection{Self-Reflection}
\label{app:self_reflection}

The choice of prompt engineering techniques has a strong impact on the quality of knowledge graph extraction ~\citep{zhu2024llms}, and different techniques have different costs in terms of tokens consumed and generated by the model.
\emph{Self-reflection} is a prompt engineering technique where the LLM generates an answer, and is then prompted to evaluate its output for correctness, clarity, or completeness, then finally generate an improved response based on that evaluation \citep{huang2023large, shinn2024reflexion, wang2022self, madaan2024self}.
We leverage self-reflection in knowledge graph extraction, and explore ways how removing self-reflection affects performance and cost.

Using larger chunk size is less costly in terms of calls to the LLM.
However, the LLM tends to extract few entities from chunks of larger size.
For example, in a sample dataset (HotPotQA,~\citeauthor{yang2018hotpotqa},~\citeyear{yang2018hotpotqa}), GPT-4 extracted almost twice as many entity references when the chunk size was 600 tokens than when it was 2400.
To address this issue, we deploy a self-reflection prompt engineering approach.
After entities are extracted from a chunk, we provide the extracted entities back to the LLM, prompting it to ``glean'' any entities that it may have missed.
This is a multi-stage process in which we first ask the LLM to assess whether all entities were extracted, using a logit bias of 100 to force a yes/no decision.
If the LLM responds that entities were missed, then a continuation indicating that ``MANY entities were missed in the last extraction'' encourages the LLM to detect these missing entities.
This approach allows us to use larger chunk sizes without a drop in quality (\autoref{fig:chunkentities}) or the forced introduction of noise.
We interate self-reflection steps up to a specified maximum number of times.

\import{}{self_reflection_figure.tex}

\section{Example Community Detection}
\label{app:graph}

\import{}{communities_figure.tex}

\raggedbottom

\pagebreak
\section{Context Window Selection}
\label{app:window}
The effect of context window size on any particular task is unclear, especially for models like \texttt{gpt-4-turbo} with a large context size of 128k tokens. Given the potential for information to be ``lost in the middle'' of longer contexts~\citep{liu-etal:2023:tacl, kuratov2024search}, we wanted to explore the effects of varying the context window size for our combinations of datasets, questions, and metrics. In particular, our goal was to determine the optimum context size for our baseline condition (\textbf{SS}) and then use this uniformly for all query-time LLM use. To that end, we tested four context window sizes: 8k, 16k, 32k and 64k. Surprisingly, the smallest context window size tested (8k) was universally better for all comparisons on comprehensiveness (average win rate of 58.1\%), while performing comparably with larger context sizes on diversity (average win rate = 52.4\%), and empowerment (average win rate = 51.3\%). Given our preference for more comprehensive and diverse answers, we therefore used a fixed context window size of 8k tokens for the final evaluation.

\raggedbottom

\pagebreak
\section{Example Answer Comparison}
\label{app:question}
\import{}{answer_table.tex}

\raggedbottom

\pagebreak
\section{System Prompts}
\label{app:sys_prompts}
\import{}{system_prompts.tex}

\raggedbottom

\pagebreak
\section{Evaluation Prompts}
\label{app:eval_prompts}
\import{}{evaluation_prompts.tex}

\raggedbottom

\pagebreak
\section{Statistical Analysis}
\label{app:stats} 
\import{}{stats_table.tex}

\raggedbottom